% TOP_PROBLEM: {"id":30007754,"problem_number":"LOCAL-30007754","title":"Wealth-tax behavioral responses","category_id":21,"category":{"id":21,"name":"economics","display_name":"Economics","description":"Open economic theory statements, published research agendas, and proposed scoped specifications; question provenance and historical priority are qualified per record.","slug":"economics","order_index":21,"created_at":"2026-10-08T00:00:00Z"},"status":"provenance_pending","external_url":null,"legacy_ids":[],"published":false,"statement":"In the explicitly specified setting: How do wealth taxes change saving, entrepreneurship, migration, avoidance, and revenue over the long run? The mathematical statement above fixes the target, assumptions and scope of this catalogue formulation.","economics":{"original_id":243,"field":"Taxation and social insurance","kind":"empirical","formalization_basis":"proposed_specification","source_status":"not_verified","priority_status":"not_established","priority_note":"An original explicit open-question statement was not verified; the cited supporting paper must not be treated as its first listing.","earliest_verified_reference":null,"current_reference":{"authors":["Arun Advani","Hannah Tarrant"],"title":"Behavioural responses to a wealth tax","year":2021,"url":"https://ifs.org.uk/journals/behavioural-responses-wealth-tax","doi":"10.1111/1475-5890.12283","locator":"Fiscal Studies 42(3–4), 509–537; IFS author abstract, paragraphs 1–2","evidence_summary":"Reviews responses and differences due to tax design, context and methodology. The accessible abstract supplies findings, but does not explicitly list the complete long-run causal target as an open problem."},"formalization_note":"Catalogue-authored scoped formalization. Population, instruments, welfare objective and identification restrictions must be instantiated before claiming a resolution; source authors are not credited with these added assumptions."},"statusQualification":"Provenance pending: an explicit published statement of this exact open question was not verified. This is a catalogue-authored candidate specification, not a certified open conjecture. Catalogue-authored scoped formalization. Population, instruments, welfare objective and identification restrictions must be instantiated before claiming a resolution; source authors are not credited with these added assumptions.","statusReviewedAt":"2026-10-08"}
% FIRST_POSTED: 2026-10-08
% ENTRY_KIND: statement_only
% !TeX program = lualatex
\documentclass[11pt]{article}
\usepackage{fontspec}
\usepackage[a4paper,margin=25mm]{geometry}
\usepackage{amsmath,amssymb,mathtools}
\usepackage{xurl}
\usepackage[hidelinks]{hyperref}
\newcommand{\sref}[2]{\hyperlink{source:#1:#2}{[S#2]}}
\setlength{\emergencystretch}{3em}
\begin{document}
% problemId:30007754
\section*{Wealth-tax behavioral responses}\label{30007754}
\subsection{Definitions and mathematical statement}
Fix a jurisdiction, implementation date, eligible population and follow-up horizon \(H\). A wealth-tax regime \(p=(\tau,b,v,e)\) specifies a rate schedule \(\tau\), exemption \(b\), asset-valuation rules \(v\), and enforcement regime \(e\). Let \(p_0\) be a specified baseline. For each person initially resident, define potential cumulative saving \(S_i^H(p)\), business investment \(I_i^H(p)\), an emigration indicator \(M_i^H(p)\in\{0,1\}\), equal to one if the person has left by \(H\), concealed assets \(A_i^H(p)\), and government receipts \(R_i^H(p)\), including receipts after emigration under the actual legal rules. Define the joint causal response
\[\Delta_H(p,p_0)=\mathbb E[(S_i^H,I_i^H,M_i^H,A_i^H,R_i^H)(p)-(S_i^H,I_i^H,M_i^H,A_i^H,R_i^H)(p_0)].\]
Specify whether intervention applies to one jurisdiction or all competing jurisdictions, since these are different counterfactuals. Identify this vector, or sharp bounds, using linked asset, ownership, residence and tax data. State the assignment or natural-experiment assumptions, treatment of anticipation, measurement-error model for private assets, and restrictions on equilibrium spillovers. Separate true wealth accumulation from reporting, valuation and relocation responses. Determine whether effects persist over prespecified horizons and transport across policy designs after conditioning on enforcement and asset composition. No assertion is made that observed taxable-wealth elasticity alone identifies this vector.

\par\textbf{Formulation and scope.} Catalogue-authored scoped formalization. Population, instruments, welfare objective and identification restrictions must be instantiated before claiming a resolution; source authors are not credited with these added assumptions.

\par\textbf{Open-question provenance.} An explicit published open-question statement was not verified. Sources below are supporting literature and must not be treated as a first listing.

\par\textbf{Historical priority.} An original explicit open-question statement was not verified; the cited supporting paper must not be treated as its first listing.

\subsection{Short English statement}
In the explicitly specified setting: How do wealth taxes change saving, entrepreneurship, migration, avoidance, and revenue over the long run? The mathematical statement above fixes the target, assumptions and scope of this catalogue formulation.

\subsection{Sources}
\begin{itemize}\raggedright
\item[\textnormal{[S1]}]\hypertarget{source:30007754:1}{} Arun Advani, Hannah Tarrant. 2021. Behavioural responses to a wealth tax. Fiscal Studies 42(3–4), 509–537; IFS author abstract, paragraphs 1–2.
\url{https://ifs.org.uk/journals/behavioural-responses-wealth-tax}.
DOI: 10.1111/1475-5890.12283.
{\small\textit{Supporting or current literature; see evidence qualification. Reviews responses and differences due to tax design, context and methodology. The accessible abstract supplies findings, but does not explicitly list the complete long-run causal target as an open problem.}}
\end{itemize}
\end{document}
